% !PRINTED 209
\weberpage{209}
Now if $S$ is any substitution of the group $\mathfrak{G}$ through which, when $a, b, c, \partial$ denote integers taken with respect to the modulus $n$, the passage

\[
x_{1,0}, x_{0,1} \quad \text{into} \quad x_{\partial,-c}, x_{-b,a}
\]

takes place, then we can apply $S$ to each of the formulae (1), (2). But by (2):

\[
x_{\mu,0} = \varphi_\mu(x_{1,0}),
\]

from which it is to be concluded that through the substitution $S$ the passage

\[
x_{\mu,0} \quad \text{into} \quad \varphi_\mu(x_{\partial,-c}) = x_{\partial\mu,\,-c\mu}
\]

takes place. Likewise one finds that through $S$ the passage

\[
x_{0,\mu'} \quad \text{into} \quad \varphi_{\mu'}(x_{-b,a}) = x_{-b\mu',\,a\mu'}
\]

takes place.

If we apply this to the equation resulting from (1)

\[
x_{\mu,\mu'} = f(x_{\mu,0}, x_{0,\mu'})
\]

then it follows that through $S$ the passage

\weq{3}{x_{\mu,\mu'} \quad \text{into} \quad x_{\partial\mu - b\mu',\,-c\mu + a\mu'}}

takes place.

From this we conclude first of all that the numbers $a, b, c, \partial$ must be so constituted that their determinant

\weq{4}{m = a\partial - bc}

has no divisor in common with $n$. For if such a divisor were present, then $\mu, \mu'$ could be so determined, without both being divisible by $n$, that

\[
\partial\mu - b\mu' \equiv 0, \quad -c\mu + a\mu' \equiv 0 \pmod{n},
\]

and several roots $x_{\mu,\mu'}$ distinct from one another would pass through $S$ into one and the same root, which is impossible. For, to begin with, neither $a$ and $b$ nor $c$ and $\partial$ can have a divisor in common with $n$, since otherwise, if $a\mu'$, $b\mu'$ or $c\mu$, $\partial\mu$ are taken divisible by $n$, all the $x_{\mu',0}$ or all the $x_{0,\mu}$ would pass into $x_{0,0}$. If one then puts $m\mu = n(ha + h'b)$, $m\mu' = n(hc + h'\partial)$, and determines $h$ and $h'$ so that $ha + h'b$ has no divisors in common with $n$, then, if $n$ and $m$ have a common divisor, $\mu, \mu'$ need not both be divisible by $n$, and not only $x_{0,0}$ but also $x_{\mu,\mu'}$ passes into $x_{0,0}$ by (3).

If we denote the permutation (3) briefly by

\[
(x_{\mu,\mu'}, x_{\nu,\nu'}),
\]

then

\[
\begin{aligned}
\nu &\equiv \partial\mu - b\mu' \\
\nu' &\equiv -c\mu + a\mu'
\end{aligned}
\pmod{n},
\]
%CONT>>
